\section*{Лекция 3. Градиентный спуск}
\subsection* {Введение}

Целью всех предыдущих занятий являлось как можно больше раскрыть процесс обучения модели. Сейчас для полной картины нам осталось разобраться с процессом подбора весов.

Ранее мы упомянали, что пока что, кажется, на ум приходит только какой"=то перебор, но согласитесь, что это довольно странно. Как минимум это странно потому, что мы не используем значения функций потерь --- неужели они нам нужны только чтобы самим посмотреть, что там у модели? Да и в конце концов, мы до сих пор не определились, в какой момент мы сможем утверждать, что текущая комбинация весов является наиболее оптимальной.

Целью этого занятия является ознакомление с процессом подбора параметров, а также знакомство с главным методом в этой сфере --- градиентный спуск.